% ======================================================================
%  frontmatter.tex —— 扉页与译校说明（不含节编号）
% ======================================================================
\begin{titlepage}
\centering
\vspace*{3.2cm}

{\sffamily\bfseries\fontsize{34}{40}\selectfont 代数几何讲义 I\par}

\vspace{0.9cm}
{\sffamily\Large 预层与层\,·\,概形\,·\,凝聚层\,·\,射影态射\par}

\vspace{2.2cm}
{\LARGE\kaishu 中译注本\par}

\vspace{2.6cm}
{\large 原著：Lin Chen\par}
\vspace{0.35cm}
{\large \textit{Notes for Algebraic Geometry 1}\par}
\vspace{0.35cm}
{\large 2025 年秋季课程讲义（底本日期：2025 年 12 月 27 日，共 171 页）\par}

\vfill

{\large 二〇二六年十一月\par}
\vspace{0.6cm}
{\small 全书 31 节\,·\,附录 A（范畴论补遗）\,·\,附录 B（术语英中对照）\par}
\end{titlepage}

\clearpage
\section*{译校说明}
\addcontentsline{toc}{section}{译校说明}
\markboth{译校说明}{译校说明}

\subsection*{一、底本}

本书是 Lin Chen（陈麟）2025 年秋季学期「代数几何 I」课程讲义的全文汉译注本。
底本为 \textit{Notes for Algebraic Geometry 1}，日期 2025 年 12 月 27 日，共 171 页，
文件 1,108,258 字节，未加密，带完整文本层。原著为一学期课程的讲义，从层论讲起，
经概形的定义与基本性质，讲到（拟）射影态射与真态射，最后以附录「Abel 范畴」收尾。
原书不含习题解答，习题（Exercise）只给陈述；本译本照录，不作解答。

\subsection*{二、页码映射}

原书的页眉即页码，且\textbf{原书页码与 PDF 页码完全一致}（偏移为 0），
故译本中凡称「原书 p.\(N\)」，均指底本 PDF 的第 \(N\) 页，读者可直接翻检。
各节在原书中的页码区间如下表；表中「节」指原书节号，译本沿用同一编号。

\begin{center}
\small
\begin{longtable}{rc@{\hspace{2.2em}}rc@{\hspace{2.2em}}rc}
\toprule
节 & 原书页 & 节 & 原书页 & 节 & 原书页\\
\midrule
\endfirsthead
\toprule
节 & 原书页 & 节 & 原书页 & 节 & 原书页\\
\midrule
\endhead
0 & 5--8 & 11 & 56--60 & 22 & 108--110\\
1 & 9--13 & 12 & 61--66 & 23 & 111--113\\
2 & 14--18 & 13 & 67--71 & 24 & 114--118\\
3 & 19--26 & 14 & 72--76 & 25 & 119--123\\
4 & 27--32 & 15 & 77 & 26 & 124--132\\
5 & 33--36 & 16 & 78--82 & 27 & 133--140\\
6 & 37--41 & 17 & 83--84 & 28 & 141--156\\
7 & 42--43 & 18 & 85--94 & 29 & 157--160\\
8 & 44--46 & 19 & 95--97 & 30 & 161--163\\
9 & 47--51 & 20 & 98--104 & 附录 A & 164--170\\
10 & 52--55 & 21 & 105--107 & 参考文献 & 171\\
\bottomrule
\end{longtable}
\end{center}

\subsection*{三、翻译原则}

\begin{enumerate}
\item \textbf{数学优先。}译文以准确传达数学内容为第一要务，不以文采取胜；
      原文的定义、命题、证明中的每一处量词、条件、方向均逐一核对。
\item \textbf{术语有定译则用定译，无定译则存原形。}凡有稳定通行汉译的术语
      （如 sheaf 层、scheme 概形、coherent 凝聚的、fiber product 纤维积），
      一律用汉译；两种以上并行译法尚无定论者（如 ample / very ample），
      保留英文并首次出现时给出译名。全书术语见书末「术语表」。
\item \textbf{不擅改原文。}原书的笔误、记号冲突、表述含混，一律保留原样，
      另以译注提示；确有把握的笔误，先判定再落笔，并在译注中说明依据。
\item \textbf{编号一致。}定义、定理、命题、引理、推论、例、注记、记号、构造、
      习题、警告共用一套计数器，按「节.小节.序号」编号，与原书编号逐条对应。
      原文中以文字方式给出的编号（如「Theorem 20.3.2」），译本改用交叉引用，
      点击可跳转，但显示的数字与原书相同。
\item \textbf{公式重排。}底本 PDF 的文本层无法直接复制公式，全部公式依原排版
      重新录入，行内公式保持行内，行间公式保持行间；交换图改用 \texttt{tikz-cd}
      重绘，节点文字与箭头方向悉依原图。
\end{enumerate}

\subsection*{四、译注体例}

译注（正文中以左侧竖线标示的灰框）只写三类内容，不写读后感，不作教材式的展开：

\begin{itemize}
\item \textbf{背景直觉}：原书一句话带过、而读者需要知道其来历或动机的概念
      （如「拟凝聚层」为何以 \(\widetilde{M}\) 为局部模型）；
\item \textbf{易错预警}：记号冲突、方向相反、条件容易漏读之处；
\item \textbf{含糊点拨}：原文省略的证明步骤、隐含的交换性、未写出的自然性。
\end{itemize}

每条解释性译注都以方括号标注依据，格式为下列四种之一：

\begin{itemize}
\item {}[来源：Stacks Project, Tag 0ABC，\url{https://stacks.math.columbia.edu/tag/0ABC}（2026-11-17 访问）。]
\item {}[来源：nLab "entry"，\url{https://ncatlab.org/nlab/show/slug}（Last revised YYYY-MM-DD；2026-11-17 访问）。]
\item {}[依据：原书 p.37 第 2 段]——指原书自身的另一处论述，用于互见；
\item {}[译者按（上下文推断，无外部依据）：……]——只用于原文自身可判定的情形。
\end{itemize}

外部依据只取三类：其一，权威教材与原始文献（Atiyah--Macdonald、Hartshorne、
Grothendieck--Dieudonné、EGA、SGA 等）；其二，Stacks Project 与 nLab 两部
公开的百科全书式文献（均标注 Tag 号或条目修订日期）；其三，本译本内部互见。
网络问答、个人笔记、内容农场一类来源，一概不取。

\subsection*{五、与原书的版式差异}

\begin{enumerate}
\item 原书每部分（Part）的标题直接排在该部分首节之上；译本将部分标题独立成页，
      以求醒目，节号与节序不变。
\item 原书的脚注在全书内连续编号（p.61 处已是第 16 号）；译本同样连续编号，脚注内容照录。
      因少数脚注随原文删除或并入译注，编号与原书可能相差一两个，正文不以脚注编号互引。
\item 译本改用系统字体排版：正文宋体，标题与强调用黑体，页眉与引文用楷体，西文与数学仍为 Latin Modern。
      全部字体子集嵌入，PDF 中的中文可以检索、复制。
\item 原书的目录只列到小节；译本目录同样列到小节，标题汉译，专名保留原文。
\item 译本删去原书的索引需求，改以书末「术语表」（英中对照，按英文字母序）代之。
\item 原书参考文献照录，条目不译；需要说明其性质者，以译注附于条目之后。
\end{enumerate}

\subsection*{六、致谢与免责}

译文与译注中的错误由译者负责，与原书作者无关。读者若发现译文失当，
请对照原书页码指出，以便订正。

\vspace{1.2em}
\noindent\hfill\begin{minipage}{0.42\linewidth}
\raggedleft\kaishu 译者\\[0.3em]
二〇二六年十一月
\end{minipage}

\clearpage
